\section{Scope of the synthesis and retained derivations}\label{sec:scope}
The central object is the attainable product--test pairing. For sparse
monomial detector spaces its rank is classified by
Theorem~\ref{thm:rank}; the observation algebra determines the future
side and the binomial tangent determines the attainable past side. The
exact causal state, streaming prediction bounds and finite-state control
construction then use that same object. At an additive collision, the
confluent pairing supplies additional scale information: its normalized
coordinates remain attainable while their physical coefficients vanish
at explicit orders. The binary ticket experiment
of Section~\ref{sec:resolved-value} joins this resolved geometry to a
common sequential payoff. The mechanical common-risk collision-bit
comparison is retained as a separate application, with its own uncensored
protocol explicitly specified.

The predecessors \cite{A1v5,A1v6} are supplied in full in the repository,
including every inherited section, proof, foundation file, response and
source manifest. Its mechanical experiment retains the whole ambient
preparation, placement failure and censoring mass. Its positive likelihood
filter, boundary-complete Bellman construction, all-stage acceptance arcs,
fixed-controller response, normalized surrogates and acceptance-cost witness
have not been removed to obtain the present theorem. They are not counted
as newly proved results of the sparse-rank argument.

In particular the fixed-controller response is not differentiation through
an optimizing policy, and the surrogate value bound does not preserve
exact sparse rank under arbitrary nonpolynomial perturbations. The older
acceptance-cost witness is not substituted for the common-risk finite-state
comparison in Section~\ref{sec:common-risk}. The full-posterior control
problem with an arbitrary hidden-parameter terminal loss remains different
from the observable future-only quotient.

The earlier finite-cylinder response and covariance derivations
\cite{A1round33} provide a useful proof discipline: the reporting space,
normalizations, derivative envelopes and limiting operation must be stated
before a limit theorem is invoked. No long-time, unweighted-seam or
many-collision assertion is imported into the present finite-horizon theorem
from their historical labels. Conversely the new finite-horizon results
are positive statements for the specified experiment and do not assert a
negative conclusion about those separate research directions.

All theorem statements in this manuscript are supported by the written
analytic proofs. The accompanying executable diagnostics test finite
instances, exact identities, sparse supports, derivative ranks, update
arithmetic and partition comparisons. They are neither a proof-assistant
certificate nor a journal-level originality or acceptance judgment.

Uniformity in this paper has a specified range. The confluent checkpoint
theorem assumes the full-future inequality \eqref{eq:full-future-condition};
the matched streaming and binary-decision laws concern the five-trial
family on the compact calibration interval $[0,1/2]$. Their constants
may depend on the fixed full-support prior and the physical protocol.
There is no claim of a prior-uniform or horizon-uniform lower constant,
no assertion of computable synthesis from approximate moment data, and
no identification of the history erasure $T_\theta$ with a one-step
mechanical sign erasure. These distinctions fix the mathematical objects
being compared; they do not remove any of the stated positive results.
